% ===== 原文 pp.45--50 =====

% [原文 p.45]

\para{电 1-形式对称性。}如上文所讨论，所有整数自旋的 Wilson 线都可以被场强的乘积完全屏蔽。因此不存在电 1-形式对称性，即
\begin{equation}
    G^{(1)} = 1 \,,
    \tag{3.108}
\end{equation}
与式 (3.88) 形成对比。

\para{'t Hooft 算符。}现在 't Hooft 算符对应于余特征标
\begin{equation}
    \phi:\ \U(1) \to \cG = \SO(3) \,.
    \tag{3.109}
\end{equation}
现在，并非所有 't Hooft 算符都能被完全屏蔽，因为我们有一个非平凡的第一同伦群
\begin{equation}
    \pi_1(\SO(3)) = \Z_2 \,.
    \tag{3.110}
\end{equation}

\begin{notebox}{$\SO(3)$ 的第一同伦群}
我们可以如下从式 (3.96) 推出式 (3.110)。考虑 $\SU(2)$ 中一条连接 $\SU(2)$ 的 $\Z_2$ 中心的两个元素的路径。由于在商去 $\Z_2$ 之后这条路径的端点被等同，这条路径下降为 $\SO(3)$ 上的一个环路。这样的环路在 $\SO(3)$ 中不可缩，因为 $\SU(2)$ 中相应的路径不可缩——这仅仅是因为不允许移动路径的端点。然而，这条路径的两倍是 $\SU(2)$ 中的一个环路，从而是可缩的。这样我们就推出了式 (3.110)。
\end{notebox}

于是，未被屏蔽的 't Hooft 算符构成的群是
\begin{equation}
    \mathsf{D}_{d-3} = \Z_2 \,.
    \tag{3.111}
\end{equation}

% [Tutor Note: p.45 式 (3.110) 可直接由覆盖空间读出：SU(2) → SO(3) 是万有覆盖，纤维为核 Z2，故 π1(SO(3)) = Z2；式 (3.96) 的路径论证只是把这一事实具体化。]

\para{磁 $(d-3)$-形式对称性。}磁 $(d-3)$-形式对称性因此是
\begin{equation}
    G^{(d-3)} = \widehat{\mathsf{D}}_{d-3} = \Z_2 \,,
    \tag{3.112}
\end{equation}
与式 (3.98) 形成对比。

% [Tutor Note: p.45 式 (3.111)--(3.112)：未被屏蔽的 't Hooft 算符群取为 π1(SO(3))，其 Pontryagin 对偶给出 Z2 的磁 (d−3)-形式对称性；对照 (3.129) 的一般公式 D_{d−3} = π1(G)。式 (3.113) 中 w2 是 Z2 值的 2-上闭链，它在闭曲面上的积分只取 0 或 1，故 exp(iπ ∫ w2) 只取正负 1，恰好实现 Z2。]

\para{拓扑面算符。}生成磁 $(d-3)$-形式对称性的拓扑面算符可以表示为
\begin{equation}
    U(\Sigma_2) = \exp\left(\ii\pi\int_{\Sigma_2} w_2\right) ,
    \tag{3.113}
\end{equation}
其中 $w_2$ 是 $\SO(3)$ 规范丛的第二 Stiefel--Whitney 类。

% [原文 p.46]

\subsubsection{一般规范群的纯 Yang--Mills 理论}

考虑一个基于单的非阿贝尔规范代数的纯 $d$ 维规范理论，即
\begin{equation}
    \fg \in \{\fsu(n), \fso(n), \fsp(n), \mathfrak{e}_6, \mathfrak{e}_7, \mathfrak{e}_8, \mathfrak{f}_4, \mathfrak{g}_2\} \,.
    \tag{3.114}
\end{equation}
规范群取为连通的李群 $\cG$，使其李代数为 $\fg$
\begin{equation}
    \mathrm{Lie}(\cG) = \fg \,.
    \tag{3.115}
\end{equation}
这样的群可以表示为
\begin{equation}
    \cG = \frac{G}{\cZ} \,,
    \tag{3.116}
\end{equation}
其中：
\begin{itemize}[leftmargin=1.4em,itemsep=1pt,topsep=1pt]
\item $G$ 是单连通群，满足
\begin{equation}
    \mathrm{Lie}(G) = \fg \,,
    \tag{3.117}
\end{equation}
它的形式是
\begin{equation}
    G \in \{\SU(n), \Spin(n), \Sp(n), E_6, E_7, E_8, F_4, G_2\}
    \tag{3.118}
\end{equation}
分别对应于 $\fg$ 所取的 (3.114) 中的形式。$G$ 单连通意味着它的第一同伦群是平凡的
\begin{equation}
    \pi_1(G) = 1 \,.
    \tag{3.119}
\end{equation}
\item $\cZ$ 是 $G$ 的中心 $Z(G)$ 的子群
\begin{equation}
    \cZ \subseteq Z(G) \,.
    \tag{3.120}
\end{equation}
\end{itemize}
(3.118) 中各群的中心汇集于表 1，其中也讨论了 $\cG$ 的各种可能性。

% [Tutor Note: p.46 式 (3.121) 是商群陈述 Z(G/Z) = Z(G)/Z；用 SU(2) 验证：G = SU(2)，Z = Z2，得 Z(SO(3)) = Z2/Z2 = 1，与 (3.149) 一致。该式成立的前提是 Z ⊆ Z(G)（见 (3.120)），故 Z 自动正规、商群有意义。]

让我们指出，规范群 $\cG$ 的中心 $Z(\cG)$ 可以表示为
\begin{equation}
    Z(\cG) = \frac{Z(G)}{\cZ} \,.
    \tag{3.121}
\end{equation}
目前我们还不引入任何物质场。

% [原文 p.47]

\begin{table}[htbp]
\centering
\small
\begin{tabular}{llll}
\toprule
单连通群 & 中心 & 表示 & 表示的荷 \\
\midrule
\multirow{3}{*}{$\SU(n)$} & \multirow{3}{*}{$\Z_n$}
  & 基础表示，$\bm{F}$ & $1 \bmod n$ \\
 & & 反对称表示，$\bm{\Lambda}^2$ & $2 \bmod n$ \\
 & & 对称表示，$\bm{S}^2$ & $2 \bmod n$ \\
\midrule
$\Sp(n)$ & $\Z_2$ & 基础表示，$\bm{F}$ & $1 \bmod 2$ \\
\midrule
\multirow{2}{*}{$\Spin(2n+1)$} & \multirow{2}{*}{$\Z_2$}
  & 矢量表示，$\bm{F}$ & $0 \bmod 2$ \\
 & & 不可约旋量表示，$\bm{S}$ & $1 \bmod 2$ \\
\midrule
\multirow{3}{*}{$\Spin(4n+2)$} & \multirow{3}{*}{$\Z_4$}
  & 矢量表示，$\bm{F}$ & $2 \bmod 4$ \\
 & & 不可约旋量表示，$\bm{S}$ & $1 \bmod 4$ \\
 & & 不可约余旋量表示，$\bm{C}$ & $3 \bmod 4$ \\
\midrule
\multirow{3}{*}{$\Spin(4n)$} & \multirow{3}{*}{$\Z_2\times\Z_2$}
  & 矢量表示，$\bm{F}$ & $(1 \bmod 2,\, 1 \bmod 2)$ \\
 & & 不可约旋量表示，$\bm{S}$ & $(1 \bmod 2,\, 0 \bmod 2)$ \\
 & & 不可约余旋量表示，$\bm{C}$ & $(0 \bmod 2,\, 1 \bmod 2)$ \\
\midrule
$E_6$ & $\Z_3$ & 基础表示，$\bm{27}$ & $1 \bmod 3$ \\
$E_7$ & $\Z_2$ & 基础表示，$\bm{56}$ & $1 \bmod 2$ \\
\bottomrule
\end{tabular}
\tabtitle{1}{这里我们列出具有非平凡中心的单连通群。对每个群，我们列出一些表示以及这些表示在中心下的荷。对于 $E_8$、$F_4$ 与 $G_2$，中心是平凡的。}
\end{table}

% [Tutor Note: p.47 表 1 自洽性检验：对 SU(2)，Λ²F（单态）与 S²F（伴随，j = 1）在 Z2 中心下的荷都是 2 mod 2 = 0，因而都是 SO(3) 的表示，与 (3.107) 只保留整数自旋一致。E8、F4、G2 的中心平凡，故按 (3.126) 其纯 Yang--Mills 理论没有电 1-形式对称性。]

\para{电 1-形式对称性。}Wilson 线算符 $\cW_R$ 由规范群 $\cG$ 的任意不可约表示 $R$ 标记。这些表示可以等同于单连通群 $G$（或等价地，规范代数 $\fg$ 的不可约表示）的、具有如下性质的不可约表示：子群
\begin{equation}
    \cZ \subseteq Z(G) \subset G
    \tag{3.122}
\end{equation}
在其上作用平凡。

由场强 $F$ 的乘积构造的非真局域算符所诱发的屏蔽会给出
\begin{equation}
    \cW_R \sim \cW_{R'}
    \tag{3.123}
\end{equation}
只要 $R$ 与 $R'$ 在中心 $Z(\cG)$ 下具有相同的荷
\begin{equation}
    q(R) = q(R') \in \widehat{Z}(\cG) \,.
    \tag{3.124}
\end{equation}
于是，未被屏蔽的 Wilson 线构成的群是
\begin{equation}
    \mathsf{D}_1 = \widehat{Z}(\cG) \,,
    \tag{3.125}
\end{equation}
这意味着 1-形式对称性群是
\begin{equation}
    G^{(1)} = Z(\cG) \,,
    \tag{3.126}
\end{equation}

% [Tutor Note: p.47 式 (3.125)--(3.126)：有限 阿贝尔群下 Pontryagin 对偶是群反射式的，故 "未屏蔽线算符的荷群 \hat{Z}(G) 再取对偶" 就回到 Z(G) 本身。这也说明为何 (3.126) 里两个 Z 都没有帽子：外层对偶已经把内层的帽子吃掉了。]

% [原文 p.48]

\para{磁 $(d-3)$-形式对称性。}'t Hooft 算符 $\cH_\phi$ 由余特征标标记
\begin{equation}
    \phi:\ \U(1) \to \cG \,.
    \tag{3.127}
\end{equation}
和之前一样，屏蔽会关联
\begin{equation}
    \cH_\phi \sim \cH_{\phi'}
    \tag{3.128}
\end{equation}
如果由 $\phi$ 与 $\phi'$ 指定的 $\cG$ 中的环路，在 $\cG$ 的群流形内部是同伦的。于是，未被屏蔽的 't Hooft 算符构成的群是
\begin{equation}
    \mathsf{D}_{d-3} = \pi_1(\cG) = \cZ \,.
    \tag{3.129}
\end{equation}
$\pi_1(\cG) = \cZ$ 这一事实来自式 (3.119)，因为商去 $\cZ$ 把 $G$ 中同伦不同的路径变成了 $\cG$ 中的环路。取 Pontryagin 对偶，我们得到 $(d-3)$-形式对称性群
\begin{equation}
    G^{(d-3)} = \widehat{\pi_1(\cG)} = \widehat{\cZ} \,.
    \tag{3.130}
\end{equation}

% [Tutor Note: p.48 式 (3.129) 的 "π1(G) = Z" 应与 (3.120) 合读：万有覆盖 G 的纤维正是 Z，故 π1(G/Z) ≅ Z（作为抽象群）。这里 Z 与 π1 都按加法群记，因此取对偶后得到 (3.130) 的 \hat{Z}。]

\para{拓扑算符。}生成电 1-形式对称性的拓扑余维-2 算符是
\begin{equation}
    U_g^{(e)} ,\qquad g \in Z(\cG) \,,
    \tag{3.131}
\end{equation}
它们可以理解为诱导出群
\begin{equation}
    G_{\min} = \frac{\cG}{Z(\cG)}
    \tag{3.132}
\end{equation}
的涡旋位形的 Dirac 弦，而这些涡旋位形不能提升为 $\cG$ 的涡旋位形。也就是说，在插入 $U_g^{(e)}(\Sigma_{d-2})$ 的周围我们有
\begin{equation}
    \int_{D^2} w_2^{\min} = g \,,
    \tag{3.133}
\end{equation}
其中 $w_2^{\min}$ 是把 $G_{\min}$ 丛提升为 $\cG$ 丛的障碍类。关于障碍类的更多细节，见式 (3.101) 之后的讨论。

% [Tutor Note: p.48 式 (3.132)--(3.134) 的障碍类要分清方向：(3.133) 的 w2^min 是 "把 G_min 丛提升为 G 丛" 的障碍，故 Dirac 弦上的磁通量取值于 Z(G) 的荷；(3.134) 的 w2 则是 "把 G 丛提升为 G 丛" 的障碍，取值于 \hat{Z} 的荷。用纯 SU(2) 验证：G_min = SU(2)/Z2 = SO(3)，g ∈ Z2 即 ∫w2 = 0 或 1，正是 SU(2) 的 Z2 Dirac 弦。]

另一方面，生成磁 $(d-3)$-形式对称性的拓扑面算符是
\begin{equation}
    U_h^{(m)}(\Sigma_2) = \int_{\Sigma_2} h(w_2) ,\qquad h \in \widehat{\cZ} \,,
    \tag{3.134}
\end{equation}
其中 $w_2$ 是把 $\cG$ 规范丛提升为 $G$ 规范丛的障碍类。

\subsubsection{包含物质场}

让我们现在在上述框架中引入物质场（它们可以具有任意的时空变换性质与任意质量）
\begin{equation}
    \phi_i ,\qquad 1 \le i \le m \,,
    \tag{3.135}
\end{equation}
它们按规范群 $\cG$ 的不可约表示 $R_i$ 变换。

% [Tutor Note: p.48 式 (3.134) 的下标 h 已取值于 \hat{Z}，但 w2 又出现在被积函数里，容易误读成 "h 乘 w2"。正确读法是 h(w2)：h 是 Z 的特征标，作用在 w2 的 Z 值上得到 U(1) 相位。]

% [原文 p.49]

\para{电 1-形式对称性。}在引入物质内容之前，未被屏蔽的 Wilson 线构成群 $\widehat{Z}(\cG)$。引入物质场会诱发额外的屏蔽。例如，表示 $R_i$ 的 Wilson 线 $\cW_{R_i}$ 可以被完全屏蔽，尽管我们有
\begin{equation}
    q(R_i) \neq 0 \in \widehat{Z}(\cG) \,.
    \tag{3.136}
\end{equation}
一般地，荷的整数线性组合\footnote{这里我们把 $\widehat{Z}(\cG)$ 视为加法群。}
\begin{equation}
    q(R_i) \in \widehat{Z}(\cG), \qquad 1 \le i \le m
    \tag{3.137}
\end{equation}
生成一个子群
\begin{equation}
    Q_{\mathrm{matter}} \subseteq \widehat{Z}(\cG) \,,
    \tag{3.138}
\end{equation}
而未被屏蔽的 Wilson 线构成的群是
\begin{equation}
    \mathsf{D}_1 = \frac{\widehat{Z}(\cG)}{Q_{\mathrm{matter}}} \,.
    \tag{3.139}
\end{equation}
电 1-形式对称性群是
\begin{equation}
    G^{(1)} = \widehat{\mathsf{D}}_1 = \widehat{\dfrac{\widehat{Z}(\cG)}{Q_{\mathrm{matter}}}} \,.
    \tag{3.140}
\end{equation}
等价地，保持拓扑的余维-2 算符恰好是对应于 $Z(\cG)$ 中那些使被完全屏蔽的 Wilson 线
\begin{equation}
    Q_{\mathrm{matter}} \subseteq \widehat{Z}(\cG)
    \tag{3.141}
\end{equation}
都不带荷的元素。因此，我们可以把 $G^{(1)}$ 表示为
\begin{equation}
    G^{(1)} = \text{$Z(\cG)$ 中使所有 $R_i$ 不变的子群} .
    \tag{3.142}
\end{equation}

% [Tutor Note: p.49 式 (3.140) 需要读成两层：内层 \hat{Z}(G) 是中心 Z(G) 的特征标群（(3.124) 的荷群），外层大帽子是对商群 D1 = \hat{Z}(G)/Q_matter 再取对偶。有限 阿贝尔群的双对偶回到原群，故 (3.140) 与 (3.142) 给出同一个 G^{(1)}。]

\para{$G^{(1)}$ 的计算。}对任意可能的规范群
\begin{equation}
    \cG \neq \Spin(4n), \qquad n \ge 2 \,,
    \tag{3.143}
\end{equation}
中心是
\begin{equation}
    Z(\cG) = \Z_g , \qquad g \ge 1 \,,
    \tag{3.144}
\end{equation}
并且我们可以把 $G^{(1)}$ 表示为
\begin{equation}
    G^{(1)} = \Z_{\gcd(g,q_1,q_2,\cdots,q_m)} \,,
    \tag{3.145}
\end{equation}
其中
\begin{equation}
    q_i \equiv q(R_i) \in \widehat{Z}(\cG) = \Z_g = \{0,1,\cdots,g-1\} \,.
    \tag{3.146}
\end{equation}

% [Tutor Note: p.49 式 (3.145) 可用例 3.10 逐一检验：SU(2) 加伴随场有 g = 2、q1 = 0，gcd(2,0) = 2，得 G^{(1)} = Z2，与 (3.148) 一致；再加基础场 q = 1，gcd(2,0,1) = 1，与 (3.151) 一致；SU(n) 加伴随场 gcd(n,0) = n，与 (3.153) 一致。]

% [原文 p.50]

\begin{examplebox}{例 3.10：$\fsu(n)$ 规范理论中的伴随与基础物质}
在 3.2 节的各个例子中，我们引用了一些关于各种规范理论电 1-形式对称性的结果。这里我们给出这些结果的推导。

首先考虑带一个伴随物质场 $\Phi$ 的 $\SU(2)$ 规范理论。由于伴随物质在中心下的荷是平凡的
\begin{equation}
    Z(\SU(2)) = \Z_2 \,,
    \tag{3.147}
\end{equation}
电 1-形式对称性群是
\begin{equation}
    G^{(1)} = Z(\SU(2)) = \Z_2 \,.
    \tag{3.148}
\end{equation}
对于 $\SO(3)$ 规范理论我们有
\begin{equation}
    Z(\SO(3)) = Z(\SU(2))/\Z_2 = 1 \,,
    \tag{3.149}
\end{equation}
因此有
\begin{equation}
    G^{(1)} = 1 \,,
    \tag{3.150}
\end{equation}
这不受加入物质场 $\Phi$ 的影响。

现在再向 $\SU(2)$ 规范理论中加入一个基础物质场 $\psi$，与伴随场 $\Phi$ 并存。基础表示在 $Z(\SU(2))$ 下有非平凡的荷，因此电 1-形式对称性群被破缺为
\begin{equation}
    G^{(1)} = 1 \,.
    \tag{3.151}
\end{equation}
现在把规范群从 $\SU(2)$ 推广到 $\SU(n)$，并只考虑取 $\SU(n)$ 伴随表示的物质场 $\Phi$。同样，伴随表示在中心下不带荷
\begin{equation}
    Z(\SU(n)) = \Z_n \,,
    \tag{3.152}
\end{equation}
因此电 1-形式对称性群是
\begin{equation}
    G^{(1)} = Z(\SU(n)) = \Z_n \,.
    \tag{3.153}
\end{equation}
\end{examplebox}

\para{磁 $(d-3)$-形式对称性。}物质场的引入不影响磁对称性，它仍由式 (3.130) 给出。

\para{拓扑算符。}生成电 1-形式对称性的拓扑余维-2 算符是
\begin{equation}
    U_g^{(e)} ,\qquad g \in G^{(1)} \,,
    \tag{3.154}
\end{equation}
它们可以理解为诱导出群
\begin{equation}
    G_{\min} = \frac{\cG}{G^{(1)}}
    \tag{3.155}
\end{equation}

% [Tutor Note: p.50 式 (3.155) 中 G_min = G/G^{(1)}：对 SU(2) 加伴随场，G = SU(2)、G^{(1)} = Z2，得 G_min = SO(3)，伴随表示正是 SO(3) 的表示；对 SU(n) 加伴随场得 PSU(n) = SU(n)/Z_n，同样支持伴随表示。这与文中 "G_min 是允许全部物质表示 R_i 的极小连通李群" 相符。]
